\begin{textsample}{NEG Dim 5 – rr – Score -1.00 – rr/rr\_ef\_46.txt}  \label{ex:f5_neg_067}
4.11.2 COMPETÊNCIAS ESPECÍFICAS DE ENSINO RELIGIOSO PARA O ENSINO FUNDAMENTAL 1.
Conhecer os aspectos estruturantes das diferentes tradições/movimentos religiosos e filosofias de vida, a partir de pressupostos científicos, filosóficos, estéticos e éticos. 2.
Compreender, valorizar e respeitar as manifestações religiosas e filosofias de vida, suas experiências e saberes, em diferentes tempos, espaços e territórios. 3.
Reconhecer e cuidar de si, do outro, da coletividade e da natureza, enquanto expressão de valor da vida. 4.
Conviver com a diversidade de crenças, pensamentos, convicções, modos de ser e viver. 5.
Analisar as relações entre as tradições religiosas e os campos da cultura, da política, da economia, da saúde, da ciência, da tecnologia e do meio ambiente. 6.
Debater, problematizar e posicionar -se frente aos discursos e práticas de \textbf{intolerância}, discriminação e violência de cunho religioso, de modo a assegurar os direitos humanos no constante exercício da cidadania e da cultura de paz.

% matched lemmas: intolerância
\end{textsample}
